% ============================================================
% cap3_synthesis_v2.tex — LLM-generated synthesis paragraphs
% Model: claude-sonnet-4-6 | Papers: top-33 (top-30 + 3 seminais)
% REVIEW CAREFULLY before inserting into cap3_slr_revised.tex
% ============================================================

% ANCHOR PARAGRAPH — seminal arc
% INSERT at start of \subsubsection{RQ2.1} or \subsection{Results/Synthesis}

The canonical origin of process mining applied to software development contexts can be traced directly to the foundational contributions of Cook and Wolf in the mid-1990s. In their seminal work, \cite{cook1995} introduced three automated process discovery methods---algorithmic grammar inference, Markov models, and neural networks---demonstrating that finite-state machine representations of software processes could be extracted automatically from event traces, substantially lowering the barrier to formal process modeling. This line of inquiry was consolidated in \cite{cook1998}, which refined the event-based discovery methodology and established the conceptual vocabulary of inferring workflow structures from historical execution data, a paradigm that remains central to the field today. Crucially, both contributions embedded stochastic reasoning---particularly Markov chain representations---within the discovery pipeline, anticipating the probabilistic extensions that would later characterise mature process mining frameworks. The translation of these principles into an explicit, named framework for software engineering contexts was accomplished by \cite{rubin2007}, who applied the ProM toolkit to source code management logs and demonstrated multi-perspective analysis spanning control-flow, informational, and organisational dimensions. Together, these three works constitute the seminal arc of IC1, establishing automated discovery from SDLC event logs as both technically feasible and analytically productive, and providing PATHCAST with its foundational methodological lineage.


% IC2+IC3 BLOCK — Stochastic techniques
% REPLACES paragraph starting 'The PDF analysis sharpens the stochastic-method picture.'
% in \subsubsection{RQ2.2: Predictive and Stochastic Techniques}

Three distinct sub-families of stochastic techniques emerge from this cluster. The first, \textbf{Markov-chain-based prediction}, encompasses models that represent software processes as discrete state transitions, including reinforcement learning agents grounded in Markov Decision Processes for pull-request outcome prediction \cite{joshi2024}, probabilistic risk quantification from repository data \cite{jeon2015}, defect removal effectiveness under code churn \cite{tian2005}, agile state-space planning and tracking \cite{tyagi2021}, release-time forecasting via generalised reliability models \cite{washizaki2015}, and Stochastic Automata Networks for Waterfall team estimation \cite{massitela2018}. The second sub-family, \textbf{stochastic Petri nets}, models concurrency and asynchrony in DevOps and CI/CD pipelines through non-Markovian specifications \cite{benmesmia2021} and sensitivity analyses that identify bottleneck stages governing delivery probability \cite{bhadra2022,bhadra2023}. The third sub-family, \textbf{Monte Carlo simulation}, drives project-level forecasting in hybrid System Dynamics models \cite{li2020}, agile risk assessment frameworks \cite{lunesu2021}, cycle-time economic analysis \cite{magennis2015}, and the PRIMAD digital-twin architecture for SDLC \cite{guinea2025}.

Despite their methodological diversity, all three sub-families share a common parametrisation strategy that constitutes a structural limitation. Transition probabilities, firing rates, and distributional parameters are invariably derived from aggregate project statistics—historical defect counts, phase-level duration averages, or summary velocity metrics—rather than from transition frequencies extracted through systematic process mining of fine-grained event logs \cite{nafreen2020,bhadra2022,li2020}. Consequently, model states remain anchored to coarse lifecycle phases such as requirements, coding, testing, and release, and cannot reflect the richer, empirically grounded activity sequences that event logs encode. This means the stochastic models are essentially hand-crafted process abstractions whose structure is assumed rather than discovered, introducing a layer of subjective design choice between observed behaviour and model. This observation constitutes \textbf{Finding F3}: \emph{stochastic forecasting models for software development are parametrised from aggregate statistics rather than from mined process structures, confining their state spaces to coarse lifecycle granularity and decoupling model topology from observed execution behaviour}.


% IC1+IC2 BLOCK — PM+stochastic integration
% APPEND after IC2+IC3 block in \subsubsection{RQ2.2}

Three studies exemplify the productive convergence of process mining and stochastic modelling at the component level. Incerto~\cite{incerto2025} demonstrated that variable-length Markov chains outperform eighteen existing stochastic conformance-checking techniques across ten of eleven benchmark datasets, establishing that higher-order sequential dependencies in event logs can be captured without sacrificing computational tractability. López-Pintado~\cite{lopezpintado2023} extended process-simulation fidelity by replacing deterministic resource-availability calendars with probabilistic counterparts discovered directly from event logs, showing that stochastic resource representations more accurately replicate empirical cycle-time distributions. Jalote~\cite{jalote2021} took a complementary perspective, modelling programmer task sequences as first-order Markov chains derived from IDE interaction logs and demonstrating that transferring the process patterns of high-productivity developers to average-productivity peers yields measurable productivity gains. Together, these contributions confirm that stochastic extensions of process models improve both explanatory and predictive validity; however, each study confines its stochastic layer to a single analytical concern—conformance, simulation, or productivity modelling—without integrating multiple stages into a unified forecasting pipeline.

The closest published work to a fully integrated multi-stage architecture is PRIMAD, proposed by Guinea-Cabrera~\cite{guinea2025}, which orchestrates pm4py-based process discovery, Akka actor-system simulation, and LightGBM supervised learning within an evolutionary Digital Twin framework for the SDLC. Preliminary results indicate promising applicability for cycle-time reduction and value-flow acceleration, suggesting that component-level integration across discovery, simulation, and prediction is technically feasible. Nevertheless, PRIMAD remains at what may be characterised as Level~2 integration: validation is reported per component rather than across the assembled pipeline, no formal inter-stage contracts or data-schema specifications govern the handoffs between discovery, simulation, and learning modules, and the architecture incorporates no dedicated Monte Carlo forecasting layer capable of propagating uncertainty estimates end-to-end. Consequently, the system does not provide calibrated probabilistic forecasts with quantified confidence intervals over future SDLC trajectories. The present review therefore concludes that the Level~3 integration niche—a formally contracted, end-to-end pipeline combining process discovery, absorbing Markov chain construction, Monte Carlo simulation, and residual correction—remains unoccupied in the extant literature.


% IC1+IC3 BLOCK — PM+forecasting (sequential)
% APPEND after IC1+IC2 block in \subsubsection{RQ2.2}

A recurring theme across this cluster is the sequential integration of process mining with predictive or simulation capabilities applied to software development contexts. \cite{gupta2017} demonstrated that combining inductive miner discovery with Markov chain-based predictive analytics enables identification of recurring failure patterns in software maintenance workflows, establishing an early precedent for stochastic augmentation of mined models. \cite{pourbafrani2021} extended this direction with SIMPT, a web-based tool that automatically generates executable process trees from event logs and supports interactive what-if simulation over activity and resource parameters, lowering the barrier to process-level experimentation. \cite{buliga2025} pushed further by coupling generative AI models with business process simulation, producing synthetic traces that are simultaneously diverse and constraint-conformant, though inter-trace dependencies such as resource contention remain partially delegated to the simulation layer. \cite{caldeira2022} demonstrated that process complexity metrics extracted from IDE and version control logs correlate meaningfully with cyclomatic complexity reduction outcomes, achieving predictive accuracies exceeding 92\%, thereby grounding process-level abstractions in measurable software quality indicators.

Despite their collective sophistication, the works in this cluster share a fundamental architectural limitation: each treats process mining and predictive modelling as consecutive, loosely coupled stages rather than components of a unified pipeline governed by formal probabilistic contracts. \cite{gupta2017} applies Markov chains after discovery without propagating uncertainty back into the process model; \cite{pourbafrani2021} simulates over fixed process trees without encoding convergence guarantees; and \cite{buliga2025} relies on simulation feasibility checks external to the generative model. Crucially, none of these systems can express, as a first-class output, the probability that a given process instance will reach completion within $k$ activity cycles — the precise forecasting primitive that distinguishes absorbing Markov chain formulations from ad hoc simulation runs. \cite{caldeira2022} similarly stops at correlation and classification, without modelling the stochastic dynamics of process trajectories. This gap motivates PATHCAST's design, wherein the mined process structure feeds directly into an absorbing Markov chain whose transient-state absorption probabilities constitute a formally defined, queryable forecasting interface.


% IC2+IC4 BLOCK — Event log sources for Markov models
% INSERT at end of \subsubsection{RQ1.2: Event Log Sources and Construction}

A complementary body of work constructs Markov models directly from GitHub repository traces, bypassing explicit process discovery in favour of manual state-space definition. Jo et al.\ \cite{jo2023} modelled developer activities in GitHub repositories as Discrete-Time Markov Chains (DTMCs), finding that commits account for 62.73\% of long-term activity, and employed probabilistic temporal logic to extract collaboration patterns relevant to project management. This approach was extended in \cite{jo2024}, where Probabilistic Model Checking with Computation Tree Logic was applied across five repositories to characterise temporal dynamics, state-transition probabilities, and key behavioural drivers, yielding transparent and explainable project forecasts. Ortu et al.\ \cite{ortu2023} complemented these structural analyses by combining Markov chains with social network analysis to identify fault-insertion and fault-fixing behavioural patterns across Apache Software Foundation projects, demonstrating that developer expertise significantly moderates defect introduction likelihood. Collectively, these studies confirm that SDLC event logs extracted from public repositories are sufficiently rich to support state-transition modelling, satisfying a foundational prerequisite for PATHCAST Stage~1. However, the event-to-state mapping adopted across these works remains unstandardised and ad hoc, representing a methodological gap that PATHCAST's process mining discovery stage aims to address systematically.


% RQ3.1 UPDATE — Integration level over top-33
% REPLACES/EXTENDS paragraph on integration levels in \subsubsection{RQ3.1}

Mapping the surveyed papers onto the L0–L3 integration scale reveals a clear stratification. The largest cluster, comprising papers that touch IC2 and IC3 without IC1 — including stochastic Petri net models of CI/CD pipelines \cite{bhadra2022,bhadra2023}, Markov-based project estimation \cite{massitela2018,tyagi2021}, and Monte Carlo risk simulation \cite{lunesu2021,magennis2015} — operates almost entirely at L1, applying stochastic modelling to software processes without any process discovery layer. Papers spanning IC1 and IC2, such as the VLMC conformance checker of \cite{incerto2025} and the probabilistic resource-calendar discoverer of \cite{lopezpintado2023}, reach L2 by coupling discovered process structure with stochastic reasoning, yet neither closes the loop toward distributional schedule forecasting. Papers combining IC1 with IC3, such as \cite{salmani2022} and \cite{caldeira2022}, use process mining descriptively to assess workflows but lack the stochastic machinery required for Monte Carlo distributional forecasts. The PRIMAD framework \cite{guinea2025} is the closest antecedent, assembling elements from all three ICs — event-log discovery, simulation, and ML prediction — but its evaluation remains component-wise, placing it at L2 rather than a fully validated end-to-end pipeline. Crucially, no paper in the reviewed corpus instantiates an L3 architecture in which process mining discovery feeds an absorbing Markov chain, whose transient distributions drive Monte Carlo forecasting, with residuals corrected by a learned model, all validated on SDLC event logs. PATHCAST is proposed as the first such L3 system.

